% from ega0-1-fr.tex
\oldpage[0\textsubscript{I}]{57}
c'est-à-dire que, pour tout homomorphisme injectif $u:N'\to N$ de
$A$-modules, l'homomorphisme
$v=1\otimes u:M\otimes_A N'\to M\otimes_A N$ est injectif. Or, $v$ est aussi
un homomorphisme de $B$-modules, et on sait que pour qu'il soit injectif, il
suffit que pour tout idéal maximal $\mathfrak{n}$ de $B$,
\[
  v_\mathfrak{n}:(M\otimes_A N')_\mathfrak{n}
  \longrightarrow(M\otimes_A N)_\mathfrak{n}
\]
soit injectif. Mais comme
\[
  (M\otimes_A N)_\mathfrak{n}
  =B_\mathfrak{n}\otimes_B(M\otimes_A N)
  =M_\mathfrak{n}\otimes_A N,
\]
$v_\mathfrak{n}$ n'est autre que l'homomorphisme
\[
  1\otimes u:M_\mathfrak{n}\otimes_A N'
  \longrightarrow M_\mathfrak{n}\otimes_A N,
\]
qui est injectif puisque $M_\mathfrak{n}$ est $A$-plat.

En particulier (en faisant $B=A$), pour qu'un $A$-module $M$ soit plat, il
faut et il suffit que $M_\mathfrak{m}$ soit $A_\mathfrak{m}$-plat pour tout
idéal maximal $\mathfrak{m}$ de $A$.
\end{env}

\begin{env}[6.3.4]
\label{0.6.3.4-fr}
Soit $M$ un $A$-module~; si $M$ est plat, et si $f\in A$ n'est pas diviseur
de $0$ dans $A$, $f$ n'annule aucun élément $\neq0$ de $M$, car
l'homomorphisme $m\mapsto f\cdot m$ s'écrit $1\otimes u$, où $u$ est la
multiplication $a\mapsto f\cdot a$ dans $A$ et $M$ est identifié à
$M\otimes_A A$~; si $u$ est injective, il en est donc de même de
$1\otimes u$ puisque $M$ est plat. En particulier, si $A$ est
\emph{intègre}, $M$ est \emph{sans torsion}.

Inversement, supposons $A$ intègre, $M$ sans torsion, et supposons que pour
tout idéal maximal $\mathfrak{m}$ de $A$, $A_\mathfrak{m}$ soit un
\emph{anneau de valuation discrète}~; alors $M$ est $A$-plat. Il suffit en
effet (6.3.3) de prouver que $M_\mathfrak{m}$ est $A_\mathfrak{m}$-plat, et
on peut donc supposer que $A$ est déjà un anneau de valuation discrète. Mais
comme $M$ est limite inductive de ses sous-modules de type fini, et que ces
derniers sont sans torsion, on peut en outre se borner au cas où $M$ est de
type fini (6.1.2). La proposition résulte dans ce cas de ce que $M$ est un
$A$-module libre.

En particulier, si $A$ est un anneau \emph{intègre},
$\varphi:A\to B$ un homomorphisme d'anneaux faisant de $B$ un $A$-module
\emph{plat} et $\neq\{0\}$, $\varphi$ est nécessairement \emph{injectif}.
Inversement, si $B$ est intègre, $A$ un sous-anneau de $B$ et si pour tout
idéal maximal $\mathfrak{m}$ de $A$, $A_\mathfrak{m}$ est un anneau de
valuation discrète, $B$ est $A$-plat.
\end{env}

\subsection{Modules fidèlement plats}
\label{subsection:0.6.4-fr}

\begin{env}[6.4.1]
\label{0.6.4.1-fr}
Pour un $A$-module $M$, les quatre propriétés suivantes sont équivalentes~:
\begin{enumerate}
  \item[$a)$] Pour qu'une suite $N'\to N\to N''$ de $A$-modules soit exacte,
    il faut et il suffit que la suite
    $M\otimes N'\to M\otimes N\to M\otimes N''$ soit exacte~;
  \item[$b)$] $M$ est plat et pour tout $A$-module $N$, la relation
    $M\otimes N=0$ entraîne $N=0$~;
  \item[$c)$] $M$ est plat et pour tout homomorphisme $v:N\to N'$ de
    $A$-modules, la relation $1_M\otimes v=0$ entraîne $v=0$, $1_M$ étant
    l'automorphisme identique de $M$~;
  \item[$d)$] $M$ est plat et pour tout idéal maximal $\mathfrak{m}$ de
    $A$, $\mathfrak{m}M\neq M$.
\end{enumerate}

Lorsque $M$ vérifie ces conditions, on dit que $M$ est un $A$-module
\emph{fidèlement plat}~; $M$ est alors nécessairement un module
\emph{fidèle}. En outre, si $u:N\to N'$ est un homomorphisme de $A$-modules,
pour que $u$ soit injectif (resp. surjectif, bijectif), il faut et il suffit
que $1\otimes u:M\otimes N\to M\otimes N'$ le soit.
\end{env}

